\section{Introduction}
\label{sec:introduction}

\subsection{Two Families of Structure-Layer Algorithms}

The redistricting platform's three-layer compositor separates plan generation
into Structure (how each bisection node is solved), WeightsOverride (what edge
weights encode), and SeedCompositor (how many seeds are tried).
The B-series has developed a rich family of Structure-layer algorithms,
all sharing a common pattern: they operate on the graph topology of the
census-tract adjacency and optimise some variant of the edge-cut objective.
\metis{} minimises raw edge cut via multilevel KL refinement~\citep{karypis1998};
\geosec{} (B.8) normalises edge cut by the isoperimetric ratio
$\sqrt{\min(i,k-i)}$ to favour balanced tree structures~\citep{b8paper};
\areasec{} (B.9) adds an area-swing constraint.
\sa{} (B.19) replaces the single \metis{} call with a simulated annealing
refinement loop to escape local optima in the edge-cut landscape~\citep{b19paper}.

All of these algorithms are, fundamentally, graph-topology methods: they
reason about which edges to cut, which vertices to separate, and how to
minimise the cost of a topological partition.

\paragraph{A second family: geometry.}
This paper introduces a different approach.
Centroidal Voronoi Districts (\cvd{}) assigns each census tract to its nearest
seed by geographic proximity and iterates seed placement to convergence.
No edge-cut objective is consulted.
The algorithm does not reason about graph topology at all during its main loop;
it reasons about \emph{distance} --- which seed is closest to each tract.
\cvd{} is in the tradition of Voronoi-based spatial partitioning rather than
graph partitioning.

The conceptual distinction matters in practice.
Edge-cut minimisation produces plans that follow existing graph structure:
METIS tends to find the ``thinest'' cuts in the adjacency graph, which often
align with roads, rivers, and county boundaries because those geographic
features create naturally low-weight edges (short shared boundaries).
\cvd{} instead fills geographic space radially from seed points.
In convex regions, \cvd{} produces roughly circular districts.
In elongated or irregular regions, it produces roughly elliptical districts
aligned with the long axis.
The resulting compactness profile is fundamentally different from the
compactness achieved by edge-cut minimisation.

\subsection{Lloyd's Algorithm and the CVD Connection}

\cvd{} is the discrete redistricting analogue of Lloyd's algorithm for
optimal quantisation~\citep{lloyd1982}.
Lloyd's algorithm, in its Euclidean form, iterates:
\begin{enumerate}
\item Assign each point to its nearest centre.
\item Move each centre to the centroid (population-weighted mean) of its cluster.
\item Repeat until convergence.
\end{enumerate}
Du, Faber and Gunzburger~\citeyear{du1999} proved convergence for Centroidal
Voronoi Tessellations (CVTs) under mild regularity conditions and surveyed their
applications in mesh generation, image compression, and spatial discretisation.

\cvd{} replaces:
\begin{itemize}
\item \textbf{Euclidean distance} with \textbf{BFS hop-count distance} on the
  census-tract adjacency graph (Phase 1, this paper);
\item \textbf{the centroid} (a possibly non-existent fractional point) with
  \textbf{the graph-distance medoid} (the actual tract minimising mean hop
  distance to all other tracts in its district).
\end{itemize}
The medoid substitution is necessary because the redistricting domain is
discrete: seeds must be actual census tracts, not fractional coordinates.
The BFS substitution avoids the need for geographic coordinate data, making
Phase~1 fully implementable with existing adjacency-list data.

\subsection{Use Cases for CVD}

Three practical use cases motivate \cvd{}:

\paragraph{Fast initial plan.}
\cvd{} runs in $O(n \times n_{\mathrm{iter}})$ time with no spanning trees,
no MCMC, and no parameter tuning beyond the iteration count.
For exploratory runs or as a starting point for \sa{} refinement, \cvd{}
provides a valid, balanced, compact plan at low cost.

\paragraph{Geometric baseline.}
The ``most geometrically natural partition'' is a defensible legal argument
in redistricting litigation.
\cvd{} provides this baseline: the plan produced by assigning each tract to
its nearest district centre, with no reference to political boundaries,
demographic composition, or incumbent addresses.
This is distinct from the \metis{} baseline, which minimises graph-topological
boundary length rather than geographic proximity.

\paragraph{Compactness diversity.}
The ensemble literature~\citep{deford2021} shows that plan quality is best
assessed by comparing a map against the space of alternatives.
Adding \cvd{} plans to the ensemble diversifies the compactness profile:
geometric proximity and edge-cut minimisation are complementary objectives
that explore different regions of the plan space.

\subsection{Main Results}

\paragraph{Algorithm.}
We define \cvd{} bisection (Section~\ref{sec:algorithm}) as a drop-in
Structure-layer replacement for \metis{}.
The algorithm has one primary tunable parameter:
\texttt{cvd\_iters} (default 20) controls the maximum number of
assignment-medoid iterations.
Seed initialisation uses k-farthest BFS selection: the first seed is derived
deterministically from the global base seed and node path; the second is the
BFS-farthest tract from the first.
A post-hoc 200-iteration boundary-swap rebalance enforces the population
balance tolerance.

\paragraph{Convergence.}
For census-tract graphs, \cvd{} typically converges in 8--14 iterations on
large states (NC, TX).
For path-like or elongated subgraphs (rural county chains), k-farthest seeding
trivially places seeds at opposite ends of the path, and convergence occurs in
1--3 steps.
The algorithm halts when seeds are unchanged between consecutive iterations.

\paragraph{Compactness profile.}
\cvd{} achieves Polsby-Popper compactness competitive with \metis{}
on NC ($k=14$) despite not optimising edge cuts.
The compactness arises naturally from the geometric proximity objective: tracts
near the centre of a district tend to be compact because they are equally
distant from all boundaries.
Section~\ref{sec:comparison} gives the full three-algorithm comparison
(\cvd{}, \metis{}, \sa{}) on NC, WI, and TX.

\paragraph{Phase~2 preview.}
Section~\ref{sec:phase2} previews Geographic \cvd{}: replacing BFS hop-count
distance with Euclidean distance on EPSG:5070 projected tract centroids.
This requires adding a \texttt{tract\_centroids} field to \texttt{LoadedGraph}
(a Phase~2 task), but is expected to yield a 5--10\% further PP improvement
for coastal and elongated states where hop-count distance diverges from
geographic distance.

\subsection{Paper Structure}

Section~\ref{sec:background} reviews Lloyd's algorithm, CVT theory, and
graph-distance medoid clustering.
Section~\ref{sec:algorithm} gives the formal \cvd{} algorithm with pseudocode,
the seed derivation, and convergence analysis.
Section~\ref{sec:comparison} presents empirical comparisons on NC, WI, and TX.
Section~\ref{sec:phase2} previews Phase~2 geographic \cvd{}.
Section~\ref{sec:conclusion} situates \cvd{} in the B-series algorithm landscape.

\subsection{Relation to Prior Work}

\cvd{} is a Structure-layer algorithm, orthogonal to the SeedCompositor
strategies (\cs{}, \be{}).
It is related to but distinct from \metis{}~\citep{karypis1998} (graph
topology vs.\ geometry) and from \sa{}~\citep{b19paper} (no warm start,
different objective).
It is the companion to \bfsgrowth{} (B.23): BFS growth from k seeds also
uses geographic proximity but grows greedily rather than iterating to
convergence.
Lloyd's algorithm in redistricting is not new as a concept, but prior
applications~\citep{du1999} use Euclidean distance and continuous coordinates;
the graph-distance medoid formulation for census-tract graphs is, to the
authors' knowledge, novel.
